% Converted from drafts/05_economics.md. Budget: 1.0 page.
\section{Economics and the Routing Confound}\label{sec:econ}

\textbf{Break-even.} At dilution rate $\epsilon$ the gateway saves
$\$1{,}000 \cdot \epsilon \cdot (1 - r)$ per \$1,000 of billed traffic, where $r = 0.115$ is the
price ratio of substitute to genuine model. Against this we set the imputed cost of enough audit
sessions to reach 80\% detection power, using measured token counts and public list prices. Both
curves are shown in Fig.~\ref{fig:F2}.

\begin{figure}[t]
\centering
\includegraphics[width=\columnwidth]{fig_break_even}
\caption{Break-even (v2, hand-written ladder). Adversary saving per \$1,000 of monthly traffic vs.\
imputed cost of an audit reaching 80\% power, by dilution rate. Hollow markers exceed the 5\%
family false-positive budget. No admissible audit exists left of $\epsilon^*$ (shaded).}
\label{fig:F2}
\end{figure}

Cost is never the obstacle. One audit session costs between \$0.0013 (KBF) and \$0.014 (FUSE), and
wherever an auditor reaches 80\% power the total is below one cent, against savings of tens of
dollars. On price alone, auditing always pays.

False positives are the obstacle. Repeating a session raises power, but it raises false alarms
too: with per-session false-positive rate $f$, $k$ sessions raise at least one false alarm with
probability $1 - (1-f)^k$. The protocol seals a family budget of 5\%, which caps each auditor at 6
to 15 sessions. Below $\epsilon = 0.10$ per-session power collapses: at $\epsilon = 0.07$ the best
auditor detects 6\% of sessions on the hand-written ladder and 1\% on the fitted one, so no
admissible number of repeats reaches 80\%. The smallest dilution an admissible audit detects is
therefore $\epsilon^* = 0.10$ on both ladders, with a bootstrap 95\% interval of [0.10, 0.10].
IRIS-lite gets there with 34\% power per session and 4 sessions (hand-written ladder) or 56\% and 2
sessions (fitted ladder), for \$0.003--\$0.006.

One caveat changes the number. IRIS-lite sets $\epsilon^*$ on both ladders, but on real recordings
its features cannot separate 20b from 120b (Section~\ref{sec:real}); its mock signal comes from
hand-written surface habits. Removing it after the seal, as a post-hoc and unsealed check, leaves
$\epsilon^* = 0.10$ on the hand-written ladder and moves it to $0.15$ on the fitted ladder, both set
by GATEOPS. We therefore claim $\epsilon^* \in [0.10, 0.15]$. A rational gateway sets $\epsilon$
just below it. At $\epsilon = 0.05$ it still saves \$44 per \$1,000, and no auditor here detects it
within a false-alarm budget a customer would accept.

\textbf{The confound.} A3 at $\epsilon = 0.10$ sends 10\% of traffic to a cheaper model. A11
spreads traffic across honest providers of the correct model. The first is fraud; the second is
ordinary load balancing, and routers do it by default. Fig.~\ref{fig:F6} and
Table~\ref{tab:confound} compare them over 200 sessions each.

\begin{figure}[t]
\centering
\includegraphics[width=\columnwidth]{fig_confound}
\caption{Share of sessions flagged, A3 (fraud, $\epsilon = 0.10$) vs.\ A11 (benign routing), v2
hand-written ladder, 200 sessions each. Top row: AUROC of A3 vs.\ A11. Timing-based scores rank
benign routing above fraud.}
\label{fig:F6}
\end{figure}

\begin{table}[t]
\centering
\caption{A3 ($\epsilon = 0.10$, fraud) vs.\ A11 (benign routing), hand-written / fitted ladder.
AUROC below 0.5 means benign routing scores as more suspicious than fraud.}
\label{tab:confound}
\small
\setlength{\tabcolsep}{4pt}
\begin{tabular}{@{}lccc@{}}
\toprule
Auditor & Flags A3 & Flags A11 & AUROC \\
\midrule
IRIS-lite & 34\% / 56\% & 0\% / 0\% & 0.96 / 0.98 \\
OTE & 4\% / 0\% & 0\% / 0\% & 0.68 / 0.48 \\
GATEOPS & 18\% / 5\% & 96\% / 100\% & 0.03 / 0.00 \\
FUSE & 2\% / 0\% & 5\% / 50\% & 0.35 / 0.01 \\
\bottomrule
\end{tabular}
\end{table}

GATEOPS finds benign routing \emph{more} anomalous than fraud (AUROC 0.00--0.03), and FUSE inherits
this: on the fitted ladder it flags half of all honest routing sessions and no fraud sessions.
Timing and metadata measure whether traffic is \emph{heterogeneous}, not whether the model is
\emph{wrong}, and honest routing is maximally heterogeneous. The text auditors separate the two arms
only because our simulated alternate provider returns near-identical text. A real alternate
provider running a different engine or quantisation would shift the text too, and the gap would
narrow.

This suggests that ``model substitution'' merges three violations that need different evidence:
\begin{itemize}
\item \emph{Identity:} the weights are wrong (A1, A2, A7). Content signals can observe this.
\item \emph{Disclosure:} routing happens but is not disclosed (A3 vs.\ A11). The two differ only in
  intent, which no client-side measurement observes.
\item \emph{Billing:} the invoice does not match the service (A8). This needs an invoice, and free
  tiers issue none.
\end{itemize}
Only the first is a detection problem for black-box auditors. The second needs a disclosure
obligation on the gateway, and the third needs verifiable metering. This reaches the recommendation
of hardware attestation in~\cite{cai2025} from the opposite direction: part of what auditors are
asked to detect is not observable from the client at all.
